\documentclass[12pt,reqno]{amsart}
\newcommand{\worksheetnumber}{3.1}
\input{../../101preamble.tex}

\DeclareMathOperator{\Sym}{Sym}

\title{Worksheet 3.1: Tensor and Symmetric Algebras}
\begin{document}
\maketitle
\thispagestyle{firstpage}

A map $\beta:V_1\times\cdots\times V_d\to U$ of vector spaces is \emph{multilinear} if it is linear in each variable when the other variables are held fixed. For $d=2$ this is exactly the notion of a bilinear map from Worksheet 2.2. For $d\geq1$, the \emph{$d$th tensor power} of a vector space $V$ is
\[
V^{\otimes d}\coloneqq\underbrace{V\otimes\cdots\otimes V}_{d},
\]
where we use the associativity isomorphisms from Worksheet 2.2 to identify the different ways of inserting parentheses. We also set $V^{\otimes0}\coloneqq\KK$. We write $v_1\otimes\cdots\otimes v_d$ for a pure tensor. As before, the pure tensors span $V^{\otimes d}$, but a general element is only a sum of pure tensors. More generally, $V_1\otimes\cdots\otimes V_d$ has the universal property you would expect: for every multilinear map $\beta:V_1\times\cdots\times V_d\to U$ there is a unique linear map $\widetilde\beta:V_1\otimes\cdots\otimes V_d\to U$ with $\widetilde\beta(v_1\otimes\cdots\otimes v_d)=\beta(v_1,\ldots,v_d)$. This lets us define operations on tensors by first checking multilinearity in the original vectors. The convention $V^{\otimes0}=\KK$ will allow an empty product of vectors to serve as a multiplicative identity, just as the empty word was the identity of the free algebra on a set in Worksheet 2.2.

It is worth distinguishing multilinearity in several variables from linearity after those variables have been identified. For instance, $(v,w)\mapsto v\otimes w$ is bilinear, but expanding $(v+w)\otimes(v+w)$ produces the cross terms $v\otimes w+w\otimes v$, so there is no reason for $v\mapsto v\otimes v$ to be linear. Notice also that a tensor power keeps track of each position separately: $v\otimes w$ and $w\otimes v$ are usually different elements of $V^{\otimes2}$. Imposing relations among the positions will produce the symmetric constructions of this worksheet, and the alternating constructions of the next.

The goal of this worksheet is to build algebras generated by a vector space. You have already seen the ingredients. In Worksheet 2.2 we constructed the tensor product so that its linear maps describe bilinear maps, and we met the free algebra $\KK\langle S\rangle$ on a set $S$, whose multiplication concatenates words. The \emph{tensor algebra} is the analogous construction starting from a vector space rather than a set: it allows arbitrary products of vectors, in which the order of the factors matters. Imposing commutativity gives the \emph{symmetric algebra}. Both constructions are characterized by universal properties, and after choosing a basis we will identify the symmetric algebra with a polynomial ring. We finish by comparing symmetric powers with duality, where the characteristic of $\KK$ plays a surprising role. We begin by proving the multilinear universal property and collecting some basic facts about tensor powers.

As in the previous worksheets, all rings have a multiplicative identity but need not be commutative, and the letter $\KK$ will denote a field. All vector spaces, linear maps, and tensor products are over $\KK$ unless we say otherwise, and vector spaces may have arbitrary dimension. As in Worksheet 2.2, a \emph{$\KK$-algebra} is a vector space with an associative, unital, $\KK$-bilinear multiplication, and a homomorphism of $\KK$-algebras is a linear map preserving multiplication and the multiplicative identity. We write $\alg_{\KK}$ for the category of $\KK$-algebras, and $\alg_{\KK}^{\mathrm{comm}}$ for its subcategory of commutative $\KK$-algebras.

\hrulefill

\begin{problems}
\item\label{s:iterated} \textbf{Iterated Tensor Products and Multilinear Maps.} Throughout this exercise let $V$ and $W$ be vector spaces.
\begin{enumerate}
\item Prove the multilinear universal property by induction on $d$. For the induction step, hold the final variable fixed, apply the universal property to the first $d-1$ variables, and check that the resulting map depends linearly on the final variable. Where does uniqueness enter?
\item Construct an isomorphism $V^{\otimes a}\otimes V^{\otimes b}\to V^{\otimes(a+b)}$ which concatenates pure tensors, including the cases $a=0$ and $b=0$. For three blocks, show that the two ways of concatenating them agree by evaluating on pure tensors.
\item Let $(e_i)_{i\in I}$ be a basis of $V$ and let $d\geq1$. Prove that the tensors $e_{i_1}\otimes\cdots\otimes e_{i_d}$ form a basis of $V^{\otimes d}$, and compute $\dim(V^{\otimes d})$ when $\dim(V)=n$ is finite. What happens when $d=0$, and what happens when $V=0$?
\item For a linear map $\phi:V\to W$, define $\phi^{\otimes d}$ on pure tensors, and set $\phi^{\otimes0}\coloneqq\Id_{\KK}$. Verify that $(\Id_V)^{\otimes d}=\Id_{V^{\otimes d}}$ and that $(\psi\phi)^{\otimes d}=\psi^{\otimes d}\phi^{\otimes d}$. For a permutation $\sigma\in S_d$, show that permuting the factors by
\[
\sigma(v_1\otimes\cdots\otimes v_d)\coloneqq v_{\sigma^{-1}(1)}\otimes\cdots\otimes v_{\sigma^{-1}(d)}
\]
defines a linear map $V^{\otimes d}\to V^{\otimes d}$, and check that these maps give an action of $S_d$ on $V^{\otimes d}$.
\end{enumerate}
\end{problems}

\hrulefill

An algebra built out of tensor powers should remember how many tensor factors each of its elements has. A \emph{graded vector space} is a vector space $A$ together with a decomposition $A=\bigoplus_{d\geq0}A_d$ into subspaces. An element of $A_d$ is called \emph{homogeneous of degree $d$}. By the definition of the direct sum from Worksheet 2.1, every element of $A$ is uniquely a finite sum of homogeneous elements, called its \emph{homogeneous components}. A \emph{graded algebra} is a graded vector space with an algebra structure satisfying $1\in A_0$ and $A_aA_b\subset A_{a+b}$ for all $a,b\geq0$. A homomorphism $\phi:A\to B$ of graded algebras is an algebra homomorphism with $\phi(A_d)\subset B_d$ for every $d$.

You already know an example. The polynomial ring $\KK[x_1,\ldots,x_n]$ is graded by total degree: the monomial $x_1^{a_1}\cdots x_n^{a_n}$ is homogeneous of degree $a_1+\cdots+a_n$. A polynomial need not be homogeneous; for instance, $1+x_1+x_1^2$ has three homogeneous components. Notice the role of the direct sum here. The polynomial ring contains only finite sums of homogeneous terms, and not arbitrary formal power series as in Worksheet 1.1.

To impose relations on a graded algebra while keeping track of degree, we need ideals which are compatible with the grading. A two-sided ideal $I$ of a graded algebra $A$ is \emph{homogeneous} if every homogeneous component of every element of $I$ also lies in $I$, or equivalently if $I=\bigoplus_d(I\cap A_d)$. In this case the quotient ring $A/I$ from Worksheet 1.1 inherits a grading, with $(A/I)_d=A_d/(I\cap A_d)$. It is natural to ask how to recognize a homogeneous ideal, and what can go wrong without homogeneity. We will see that it is enough to check the generators of an ideal, and that a non-homogeneous relation can destroy the grading entirely. We check these statements in the following exercise.

\hrulefill

\begin{problems}
\item\label{s:grading} \textbf{Gradings and Homogeneous Ideals.} Throughout parts (b) and (c) let $A=\bigoplus_{d\geq0}A_d$ be a graded algebra.
\begin{enumerate}
\item Verify that total degree makes $\KK[x_1,\ldots,x_n]$ a graded algebra. Find the homogeneous components of $(1+x+y)(x-y)$ in $\KK[x,y]$. Does your answer depend on the characteristic of $\KK$?
\item Prove that a two-sided ideal generated by homogeneous elements is homogeneous. Hint: write an element of the ideal as a finite sum of terms $arb$, with $r$ a homogeneous generator, and decompose $a$ and $b$ into homogeneous components.
\item Let $I\subset A$ be a homogeneous ideal. Prove that the natural map $\bigoplus_dA_d/(I\cap A_d)\to A/I$ is an isomorphism of vector spaces. Show that multiplication respects the resulting grading, and that the quotient map $A\to A/I$ preserves degree.
\item Show that $\langle x-1\rangle\subset\KK[x]$ is not homogeneous. Explain why the relation $x=1$ prevents the usual degrees of $x$ and $1$ from giving a grading on the quotient. Compare this with the ideal $\langle x^2\rangle$ and the dual numbers from Worksheet 1.1.
\end{enumerate}
\end{problems}

\hrulefill

The \emph{tensor algebra} of $V$ is the graded vector space
\[
T(V)\coloneqq\bigoplus_{d\geq0}V^{\otimes d}=\KK\oplus V\oplus(V\otimes V)\oplus\cdots.
\]
Multiplication concatenates pure tensors and extends bilinearly. More precisely, on $V^{\otimes a}\times V^{\otimes b}$ it is induced by the isomorphism of Problem~\ref{s:iterated}(b), and multiplication by an element of $V^{\otimes0}=\KK$ is scalar multiplication. The element $1\in\KK=V^{\otimes0}$ is the multiplicative identity, and there is a linear inclusion $\iota_V:V\to T(V)$ onto degree one. We will often write $v_1\cdots v_d$ for the tensor $v_1\otimes\cdots\otimes v_d$, viewed as a product in $T(V)$.

Notice that $V$ already has an addition and a scalar multiplication. To generate an algebra from $V$ we want to keep these operations, while allowing arbitrary products of vectors. Thus $\iota_V(v+w)=\iota_V(v)+\iota_V(w)$, but there is no reason for $\iota_V(v)\iota_V(w)$ to equal $\iota_V(w)\iota_V(v)$. For a one-dimensional space $V=\KK e$, an element of $T(V)$ has the form $a_0+a_1e+a_2(e\otimes e)+\cdots+a_me^{\otimes m}$, and concatenation adds the number of copies of $e$, much as multiplication adds exponents in a polynomial ring. With two basis vectors, however, degree two already has four tensor basis elements: $e_1e_1$, $e_1e_2$, $e_2e_1$, and $e_2e_2$. In particular, the two mixed products remain distinct. This is the difference between polynomials in one variable and words in several letters.

The following universal property says that the tensor algebra imposes no relations beyond those already present in $V$.

\begin{theorem}[Universal Property of the Tensor Algebra]
For every $\KK$-algebra $A$ and every linear map $\phi:V\to A$, there is a unique algebra homomorphism $\widetilde\phi:T(V)\to A$ such that $\widetilde\phi\iota_V=\phi$. That is, there is a unique dashed arrow making the following diagram commute:
\[
\begin{tikzcd}[row sep=4em,column sep=4em]
V \arrow[r,"\phi"] \arrow[d,"\iota_V"'] & A \\
T(V) \arrow[ur,dashed,"\widetilde\phi"'] &
\end{tikzcd}
\]
\end{theorem}

Notice that $A$ need not be graded, and that $\phi$ may take values anywhere in $A$. If $A$ is graded and $\phi(V)\subset A_1$, we will see that the extension preserves degrees. In degree zero the extension is forced on us by the requirement that it preserve the multiplicative identity. We prove the theorem in the following exercise.

\hrulefill

\begin{problems}
\item\label{s:tensoralgebra} \textbf{The Tensor Algebra and Its Universal Property.} Throughout this exercise let $V$ be a vector space.
\begin{enumerate}
\item Verify that the proposed multiplication makes $T(V)$ an associative, unital, graded algebra. Explain why the product of two elements still has only finitely many nonzero homogeneous components.
\item Let $\phi:V\to A$ be a linear map to a $\KK$-algebra. Prove that $(v_1,\ldots,v_d)\mapsto\phi(v_1)\cdots\phi(v_d)$ is multilinear. Construct a linear map on each degree, specifying the map $\KK\to A$ in degree zero, and combine them to define $\widetilde\phi:T(V)\to A$.
\item Check that $\widetilde\phi$ preserves multiplication and the multiplicative identity. Prove uniqueness using the fact that products of elements of degree one span every positive degree. This proves the theorem.
\item If $A$ is graded and $\phi(V)\subset A_1$, prove that $\widetilde\phi$ preserves degrees. Explain why this conclusion can fail for a general linear map $\phi:V\to A$.
\item Identify $T(0)$ and $T(\KK e)$ as algebras. If $e_1,e_2\in V$ are linearly independent, use the tensor basis to show that $e_1e_2-e_2e_1\ne0$ in $T(V)$. For which finite-dimensional $V$ is $T(V)$ commutative?
\end{enumerate}
\end{problems}

\hrulefill

The universal property also supplies maps between tensor algebras. Given a linear map $\phi:V\to W$, the composite $\iota_W\phi:V\to T(W)$ extends to an algebra homomorphism $T(\phi):T(V)\to T(W)$, which is $\phi^{\otimes d}$ in degree $d$. As with products and cokernels in Worksheet 2.1, uniqueness of the extension makes the functor conditions nearly automatic.

We can also phrase the universal property as an adjunction, in the sense of Worksheet 2.2. Write $\cU:\alg_{\KK}\to\Vect_{\KK}$ for the forgetful functor, which remembers only the underlying vector space of an algebra. The theorem gives a bijection
\[
\Theta_{V,A}:\Hom_{\alg_{\KK}}(T(V),A)\longrightarrow\Hom_{\KK}(V,\cU(A)),\qquad F\longmapsto F\iota_V.
\]
Once we check naturality, this says that $T\dashv\cU$. Compare this with the free algebra $\KK\langle S\rangle$ from Worksheet 2.2, which is left adjoint to the forgetful functor $\alg_{\KK}\to\Set$: there a function on the letters extends freely, whereas the input to the tensor algebra must be a linear map. Taking the free vector space on $S$ relates the two constructions.

It is worth pausing on this difference. Specifying the image of a vector $v\in V$ already prescribes the image of every scalar multiple of $v$, and specifying the images of $v$ and $w$ prescribes the image of $v+w$. So the vectors of $V$ are not independent generators of $T(V)$. A basis resolves this: every vector is uniquely a linear combination of basis vectors, and the images of the basis vectors really can be chosen arbitrarily. This is why a basis gives a description of $T(V)$ in terms of words, even though neither the definition of $T(V)$ nor its universal property requires one. We make these comparisons precise in the following exercise.

\hrulefill

\begin{problems}
\item\label{s:tensorfunctor} \textbf{Functoriality and the Free Algebra Adjunction.} Throughout this exercise let $V$ and $W$ be vector spaces.
\begin{enumerate}
\item Construct $T(\phi)$ for a linear map $\phi:V\to W$. Prove that $T(\Id_V)=\Id_{T(V)}$ and $T(\psi\phi)=T(\psi)T(\phi)$ by comparing the maps on scalars and in degree one.
\item Verify that $F\mapsto F\iota_V$ and $\phi\mapsto\widetilde\phi$ are inverse to one another, so that $\Theta_{V,A}$ is a bijection. Is $\Hom_{\alg_{\KK}}(T(V),A)$ a vector space under pointwise operations? Explain why this adjunction is only a bijection of sets, unlike the tensor--Hom correspondence of Worksheet 2.2.
\item Let $a:V'\to V$ be linear, let $b:A\to A'$ be an algebra homomorphism, and let $F:T(V)\to A$ be an algebra homomorphism. Prove the naturality equation
\[
\Theta_{V',A'}\bigl(bFT(a)\bigr)=\cU(b)\,\Theta_{V,A}(F)\,a.
\]
In which variable is the correspondence contravariant, and in which is it covariant?
\item Let $S$ be a set. Construct an algebra isomorphism $T(\KK^{\langle S\rangle})\cong\KK\langle S\rangle$ carrying $e_s$ to the one-letter word $s$. Construct the maps in both directions, and use the fact that both algebras are generated by their degree one elements to show that they are inverse. For a basis $B$ of $V$, deduce that $T(V)\cong\KK\langle B\rangle$.
\item Compare $T(V)$ with $\KK\langle|V|\rangle$, where $|V|$ is the underlying set of $V$ as in Worksheet 2.2. Describe the surjective algebra homomorphism $\KK\langle|V|\rangle\to T(V)$, and exhibit elements of its kernel coming from addition and scalar multiplication in $V$. What happens to the letter indexed by the zero vector?
\end{enumerate}
\end{problems}

\hrulefill

The tensor algebra permits us to send basis vectors to arbitrary elements of an algebra, including elements which do not commute. What changes if we are only interested in commutative target algebras? Every algebra homomorphism into a commutative algebra must send $vw$ and $wv$ to the same element. We can impose these equalities in the source using a quotient, exactly as for quotient rings and modules. Let $J\subset T(V)$ be the two-sided ideal generated by
\[
v\otimes w-w\otimes v\qquad(v,w\in V).
\]
The \emph{symmetric algebra} of $V$ is $\Sym(V)\coloneqq T(V)/J$. Since its relations are homogeneous of degree two, this is a graded algebra. Its degree $d$ part is the \emph{$d$th symmetric power} $\Sym^d(V)$; in particular, $\Sym^0(V)=\KK$. We write $v_1\cdots v_d$ for the image of a pure tensor. Notice that this notation now allows us to reorder the factors, which was not allowed in $T(V)$.

The quotient is commutative, because its degree one generators commute. Conversely, a linear map from $V$ into a commutative algebra extends to $T(V)$, and the extension sends each generator of $J$ to zero. The universal property of the quotient therefore gives the following result. Notice that the target is still an ordinary algebra; no grading is required.

\begin{theorem}[Universal Property of the Symmetric Algebra]
For every commutative $\KK$-algebra $A$, restriction to $V$ gives a bijection
\[
\Hom_{\alg_{\KK}^{\mathrm{comm}}}(\Sym(V),A)\cong\Hom_{\KK}(V,\cU(A)).
\]
Equivalently, every linear map $V\to A$ extends uniquely to an algebra homomorphism $\Sym(V)\to A$.
\end{theorem}

Writing $q:T(V)\to\Sym(V)$ for the quotient map, the two successive extensions fit into the diagram
\[
\begin{tikzcd}[row sep=4em,column sep=4em]
V \arrow[r,"\iota_V"] \arrow[dr,"\phi"'] & T(V) \arrow[d,"\widetilde\phi"] \arrow[r,"q"] & \Sym(V) \arrow[dl,dashed,"\overline\phi"] \\
& A &
\end{tikzcd}
\]
The first step uses the linearity of $\phi$ to construct an algebra homomorphism on tensors, and the second uses the commutativity of $A$ to ensure that this homomorphism descends. Thus the two universal properties separate two requirements on the target: it must have a multiplication in which products of vectors can be evaluated, and those products must not depend on the order of the vectors.

There is also a description one degree at a time. For $d\geq1$, a multilinear map $V^d\to U$ is \emph{symmetric} if permuting its arguments does not change its value. We will see that linear maps out of $\Sym^d(V)$ correspond to symmetric multilinear maps. To justify this we need to relate the degree $d$ part of $J$ to permutations of tensor factors, and not merely to know that the quotient algebra is commutative. We prove the theorem, and this description of symmetric powers, in the following exercise.

\hrulefill

\begin{problems}
\item\label{s:symmetric} \textbf{Construction and Universal Properties of the Symmetric Algebra.} Throughout this exercise let $V$ be a vector space and let $J\subset T(V)$ be the ideal defined above.
\begin{enumerate}
\item Prove that $J$ is homogeneous, and that $J\cap V^{\otimes0}=J\cap V^{\otimes1}=0$. Deduce that the map $V\to\Sym(V)$ identifies $V$ with $\Sym^1(V)$. Prove that arbitrary products of elements of degree one commute in $\Sym(V)$, and deduce that $\Sym(V)$ is commutative.
\item Let $A$ be a commutative algebra. Extend a linear map $V\to A$ to $T(V)$, and prove that the extension kills $J$. Use the universal property of the quotient to factor it through $\Sym(V)$, and prove uniqueness. This proves the theorem.
\item For $d\geq2$, show that $J\cap V^{\otimes d}$ is spanned by the differences between pure tensors obtained by interchanging two adjacent factors. Using the fact that adjacent transpositions generate $S_d$, show that $J\cap V^{\otimes d}$ equals
\[
N_d\coloneqq\operatorname{span}\left\{\sigma t-t \;\; \big| \;\; t\in V^{\otimes d},\ \sigma\in S_d\right\}.
\]
Deduce that $\Sym^d(V)=V^{\otimes d}/N_d$. Why does this also hold for $d=0$ and $d=1$ with $N_d=0$?
\item For $d\geq1$, use part (c) and Problem~\ref{s:iterated} to prove that linear maps $\Sym^d(V)\to U$ correspond bijectively to symmetric multilinear maps $V^d\to U$. Which symmetric multilinear map corresponds to the inclusion $\Sym^d(V)\hookrightarrow\Sym(V)$?
\item Let $\phi:V\to W$ be linear, and write $J_V\subset T(V)$ and $J_W\subset T(W)$ for the ideals defining the two symmetric algebras. Show that $T(\phi)(J_V)\subset J_W$. Construct $\Sym(\phi)$, state its value on $v_1\cdots v_d$, and prove the functor conditions. Verify the naturality of the bijection in the theorem, as in Problem~\ref{s:tensorfunctor}(c), to obtain an adjunction $\Sym\dashv\cU$ between $\Vect_{\KK}$ and $\alg_{\KK}^{\mathrm{comm}}$.
\end{enumerate}
\end{problems}

\hrulefill

Suppose $V$ is finite dimensional, with basis $e_1,\ldots,e_n$. Since the images of the $e_i$ commute in $\Sym(V)$, every product of basis vectors can be written as a monomial $e_1^{a_1}\cdots e_n^{a_n}$, and these monomials span. We will see that the universal property forces them to be linearly independent as well, so that $\Sym(V)$ is isomorphic to the polynomial ring $\KK[x_1,\ldots,x_n]$ by sending $e_i$ to $x_i$. Notice that this identification depends on the chosen basis, and is not part of the definition of $\Sym(V)$.

As in Worksheet 1.1, these are formal polynomials. Over a finite field distinct polynomials can define the same function, as you saw for $x^2-x$ over $\FF_2$, so we cannot prove linear independence by evaluating polynomial functions. Instead we will construct mutually inverse algebra homomorphisms. Once we have a monomial basis, computing the dimension of a symmetric power becomes a counting problem: we count the exponent lists with total sum $d$.

It is worth being careful about what the variables mean. The symbols $x_i$ correspond to the vectors $e_i$ themselves, and not to coordinate functions on $V$. The coordinate functionals $e_i^{\vee}$ live in the dual space, whose symmetric algebra is $\Sym(V^{\vee})$. Both algebras become polynomial rings once a basis is chosen, but comparing them uses that choice, and we will return to their relationship at the end of this worksheet. We compute bases of symmetric powers in the following exercise.

\hrulefill

\begin{problems}
\item\label{s:polynomial} \textbf{Polynomial Rings and Bases of Symmetric Powers.} Throughout this exercise let $e_1,\ldots,e_n$ be a basis of $V$.
\begin{enumerate}
\item Use the universal property of the symmetric algebra to construct an algebra homomorphism $\Sym(V)\to\KK[x_1,\ldots,x_n]$ with $e_i\mapsto x_i$, and use substitution of commuting elements to construct a homomorphism in the other direction. Check that both composites fix scalars and generators, and deduce that $\Sym(V)\cong\KK[x_1,\ldots,x_n]$ as graded algebras.
\item Prove that the monomials $e_1^{a_1}\cdots e_n^{a_n}$ with $a_i\geq0$ and $a_1+\cdots+a_n=d$ form a basis of $\Sym^d(V)$. Deduce that the quotient map $V^{\otimes d}\to\Sym^d(V)$ need not be injective. When is it injective?
\item For $n\geq1$, represent a monomial of degree $d$ by a string of $d$ identical marks separated into $n$ groups by $n-1$ dividers, and prove that
\[
\dim(\Sym^d(V))=\binom{n+d-1}{d}.
\]
Check the case $d=0$ explicitly. What is $\dim(\Sym^d(V))$ when $n=0$, for $d=0$ and for $d>0$?
\item List bases of $\Sym^2(\KK^2)$ and $\Sym^3(\KK^2)$. Compare their dimensions with those of the corresponding tensor powers, and describe which tensor basis elements have been identified.
\end{enumerate}
\end{problems}

\hrulefill

Functoriality now becomes substitution into polynomial expressions. The algebra homomorphism $\Sym(\phi)$ preserves degree, so it restricts to linear maps $\Sym^d(\phi):\Sym^d(V)\to\Sym^d(W)$. Its definition does not require bases, but monomial bases let us compute it by expanding products of the images of basis vectors. Coefficients such as $2$ appear because different ordered tensors become the same monomial in the quotient. This is our first glimpse of the characteristic of $\KK$ affecting a multilinear construction. We compute some examples in the following exercise.

\hrulefill

\begin{problems}
\item\label{s:computation} \textbf{Induced Maps on Symmetric Powers.} Throughout this exercise let $V$ and $W$ be vector spaces over $\QQ$ with bases $e_1,e_2$ and $f_1,f_2$, and let $\phi:V\to W$ be given by $\phi(e_1)=f_1+f_2$ and $\phi(e_2)=f_2$.
\begin{enumerate}
\item Compute $\Sym^2(\phi)$ on the basis $(e_1^2,e_1e_2,e_2^2)$, and write its matrix using the corresponding ordered basis $(f_1^2,f_1f_2,f_2^2)$. Compute the matrix of $\Sym^3(\phi)$ in the monomial bases of degree three.
\item Find $\phi^{-1}$, and use it to construct the inverse of $\Sym^2(\phi)$. Verify that it agrees with the inverse of your matrix. Which functor condition guarantees this agreement?
\item Compute $\phi^{\otimes2}(e_1\otimes e_1)$ before passing to the symmetric quotient. Identify the terms which combine, and verify that the square relating $\phi^{\otimes2}$ and $\Sym^2(\phi)$ to the two quotient maps commutes.
\item Repeat the degree two calculation over a field of characteristic $2$, using the same formulas for $\phi$. Which coefficient changes? Is the induced map still invertible? Justify your answer using the original linear map.
\end{enumerate}
\end{problems}

\hrulefill

The symmetric algebra also interacts well with direct sums. To describe how, recall from Worksheet 2.2 that the tensor product of two vector spaces is again a vector space. If $A$ and $B$ are $\KK$-algebras, then $A\otimes B$ is a $\KK$-algebra with multiplication and identity
\[
(a\otimes b)(a'\otimes b')\coloneqq aa'\otimes bb',\qquad 1\coloneqq1\otimes1.
\]
If $A$ and $B$ are graded, we give $A\otimes B$ the total grading $(A\otimes B)_d\coloneqq\bigoplus_{a+b=d}A_a\otimes B_b$. Notice that no sign appears when factors pass one another; we will meet a signed version in the next worksheet. In particular, the tensor product of two commutative algebras is commutative.

The inclusions of $V$ and $W$ into $V\oplus W$ induce homomorphisms out of their symmetric algebras, and multiplying their images suggests an isomorphism
\[
\Sym(V)\otimes\Sym(W)\xrightarrow{\ \sim\ }\Sym(V\oplus W).
\]
After choosing bases this is not surprising: the polynomial variables are divided into those coming from $V$ and those coming from $W$. We will instead prove it using universal properties, so that its compatibility with linear maps is also apparent.

An element of degree $d$ in $\Sym(V\oplus W)$ can have some factors from $V$ and the remaining factors from $W$. Grouping the factors according to their source suggests one summand for each pair $a+b=d$. In degree two, for instance, there are three possibilities: two factors from $V$, one from each space, or two from $W$. The algebra isomorphism will show that this grouping is a direct sum decomposition, and not merely a way of writing elements which might introduce additional relations. We prove this in the following exercise.

\hrulefill

\begin{problems}
\item\label{s:directsum} \textbf{Symmetric Algebras of Direct Sums.} Throughout this exercise let $V$ and $W$ be vector spaces.
\begin{enumerate}
\item Let $A$ and $B$ be $\KK$-algebras. Use the multilinear universal property to show that the multiplication on $A\otimes B$ is well defined. Verify associativity and the identity axiom, and show that $A\otimes B$ is commutative when $A$ and $B$ are. When $A$ and $B$ are graded, verify that the total grading makes $A\otimes B$ a graded algebra, and explain why every tensor has only finitely many homogeneous components.
\item Let $i_V$ and $i_W$ be the coordinate inclusions into $V\oplus W$. Construct an algebra homomorphism $\alpha:\Sym(V)\otimes\Sym(W)\to\Sym(V\oplus W)$ with $\alpha(a\otimes b)=\Sym(i_V)(a)\Sym(i_W)(b)$. Where do you use commutativity to prove that $\alpha$ is multiplicative?
\item Use the linear map $(v,w)\mapsto v\otimes1+1\otimes w$ to construct an algebra homomorphism $\beta$ in the other direction. Show that $\alpha\beta$ and $\beta\alpha$ fix scalars and algebra generators, and hence are identity maps.
\item Prove that $\alpha$ and $\beta$ preserve degree, and deduce that there is a natural isomorphism
\[
\Sym^d(V\oplus W)\cong\bigoplus_{a+b=d}\Sym^a(V)\otimes\Sym^b(W).
\]
Write both directions on products of vectors, and verify compatibility with linear maps $V\to V'$ and $W\to W'$.
\item Suppose $\dim(V)=m$ and $\dim(W)=n$ are finite. Take dimensions in part (d) to obtain an identity of binomial coefficients. When $\dim(V)=\dim(W)=1$, verify the decomposition in degree two directly using monomial bases.
\end{enumerate}
\end{problems}

\hrulefill

We finish by comparing symmetric powers with duality. For the rest of this worksheet, suppose $V$ is finite dimensional and let $d\geq0$. Evaluation of functionals gives a pairing
\[
(\lambda_1\otimes\cdots\otimes\lambda_d,\ v_1\otimes\cdots\otimes v_d)\longmapsto\prod_{i=1}^d\lambda_i(v_i),
\]
which identifies $(V^{\vee})^{\otimes d}$ with $(V^{\otimes d})^{\vee}$, as you can check in tensor bases. This pairing does not descend to symmetric powers in both variables, since permuting only one of the two lists can change its value. Summing over permutations gives a pairing which does descend:
\[
\left\langle\lambda_1\cdots\lambda_d,\ v_1\cdots v_d\right\rangle\coloneqq\sum_{\sigma\in S_d}\prod_{i=1}^d\lambda_i(v_{\sigma(i)}).
\]
For $d=0$ the formula is multiplication in $\KK$, since there is exactly one permutation of the empty set and an empty product equals $1$.

A bilinear pairing $X\times Y\to\KK$ between finite-dimensional spaces is \emph{perfect} if the induced maps $X\to Y^{\vee}$ and $Y\to X^{\vee}$ are both isomorphisms. It is natural to ask whether the pairing above is perfect, so that it identifies $\Sym^d(V^{\vee})$ with the dual of $\Sym^d(V)$. Repeated factors suggest that factorials should appear. Interchanging two copies of the same functional does not produce a different tensor, but both permutations still occur in the sum, and a list with several repeated functionals has several such groups of interchangeable positions. We will see that the pairing is perfect when $d!$ is invertible in $\KK$, but that it can fail badly in positive characteristic. Notice that such a failure concerns this particular pairing: the two symmetric powers still have the same dimension, by Problem~\ref{s:polynomial}. We compute these multiplicities in the following exercise.

\hrulefill

\begin{problems}
\item\label{s:pairing} \textbf{The Symmetric Pairing and Factorials.} Throughout this exercise let $V$ be finite dimensional.
\begin{enumerate}
\item Prove that the displayed formula is multilinear in the $\lambda_i$ and the $v_i$, and unchanged by a permutation of either list. Use Problem~\ref{s:symmetric} to obtain a bilinear pairing $\Sym^d(V^{\vee})\times\Sym^d(V)\to\KK$.
\item Let $e_1,\ldots,e_n$ be a basis of $V$ with dual basis $e_1^{\vee},\ldots,e_n^{\vee}$. For exponent lists $a$ and $b$ of total degree $d$, compute the pairing of $(e_1^{\vee})^{a_1}\cdots(e_n^{\vee})^{a_n}$ with $e_1^{b_1}\cdots e_n^{b_n}$. By counting the contributing permutations, show that it is $0$ when $a\ne b$ and $a_1!\cdots a_n!$ when $a=b$.
\item Deduce that the pairing is perfect if $d!$ is invertible in $\KK$. Suppose $\dim(V)=1$ and $\operatorname{char}(\KK)=p>0$. Compute the pairing in degree $p$, and show that it is identically zero, although both spaces are nonzero.
\item Let $\phi:V\to W$ be a linear map between finite-dimensional spaces. For $\lambda\in\Sym^d(W^{\vee})$ and $z\in\Sym^d(V)$, verify that
\[
\left\langle\Sym^d(\phi^{\vee})(\lambda),z\right\rangle=\left\langle\lambda,\Sym^d(\phi)(z)\right\rangle.
\]
Explain why, when $d!$ is invertible, the resulting isomorphism $\Sym^d(V^{\vee})\cong\Sym^d(V)^{\vee}$ is natural. Hint: draw the corresponding naturality square, as in Worksheet 2.1.
\end{enumerate}
\end{problems}

\hrulefill

Is there a description of $\Sym^d(V)^{\vee}$ which works in every characteristic? The answer distinguishes two ways of making a group action disappear. If a group $G$ acts linearly on a vector space $E$, its \emph{invariants} and \emph{coinvariants} are
\[
E^G\coloneqq\left\{e\in E \;\; \big| \;\; ge=e\text{ for all }g\in G\right\}
\qquad\text{and}\qquad
E_G\coloneqq E/\operatorname{span}\left\{ge-e \;\; \big| \;\; g\in G,\ e\in E\right\}.
\]
Taking invariants keeps the fixed vectors inside the original space, while taking coinvariants identifies vectors related by the action. By Problem~\ref{s:symmetric}(c), $\Sym^d(V)$ is exactly the space of coinvariants $(V^{\otimes d})_{S_d}$.

The \emph{$d$th divided power space} is the space of invariants
\[
\Gamma^d(V^{\vee})\coloneqq\bigl((V^{\vee})^{\otimes d}\bigr)^{S_d},
\]
using the permutation action from Problem~\ref{s:iterated}(d). We will need only these vector spaces and their induced maps here. Recall from Worksheet 2.1 that a functional on a quotient is the same thing as a functional on the original space which kills the relations. Under the tensor pairing, killing the permutation relations turns out to be the same as being invariant. Consequently, there is an isomorphism
\[
\Gamma^d(V^{\vee})\xrightarrow{\ \sim\ }\Sym^d(V)^{\vee},
\]
defined without choices, in every characteristic. The symmetric pairing of Problem~\ref{s:pairing} factors through the map from coinvariants to invariants which sums over permutations, and it is in this intermediate map that the factorials occur. We prove these statements in the following exercise.

\hrulefill

\begin{problems}
\item\label{s:divided} \textbf{Invariant Tensors and the Dual of a Symmetric Power.} Throughout this exercise let $V$ be finite dimensional, and write $N_d\subset V^{\otimes d}$ as in Problem~\ref{s:symmetric}(c).
\begin{enumerate}
\item Prove that the tensor pairing between $(V^{\vee})^{\otimes d}$ and $V^{\otimes d}$ is perfect, and that $\langle\sigma\xi,t\rangle=\langle\xi,\sigma^{-1}t\rangle$ for all $\sigma\in S_d$. Deduce that an invariant tensor $\xi$ defines a functional on $V^{\otimes d}/N_d$ which does not depend on the chosen representative.
\item Conversely, pull a functional on $\Sym^d(V)$ back to $V^{\otimes d}$, and represent it by a unique $\xi\in(V^{\vee})^{\otimes d}$. Since this functional kills every $\sigma t-t$, prove that $\xi$ is invariant. Deduce the asserted isomorphism, including in degree zero.
\item For an exponent list $a$ of total degree $d$, let $\gamma_a$ be the sum of the \emph{distinct} tensors $e_{i_1}^{\vee}\otimes\cdots\otimes e_{i_d}^{\vee}$ containing exactly $a_i$ copies of $e_i^{\vee}$ for each $i$. Prove that the $\gamma_a$ form a basis of $\Gamma^d(V^{\vee})$ by comparing coefficients on each orbit of $S_d$. Evaluate $\gamma_a$ on the monomials $e_1^{b_1}\cdots e_n^{b_n}$, and identify the dual basis.
\item Define $\mathcal N:\Sym^d(V^{\vee})\to\Gamma^d(V^{\vee})$ by $\mathcal N([\xi])\coloneqq\sum_{\sigma\in S_d}\sigma\xi$. Prove that $\mathcal N$ is well defined, compute it on each monomial in the $e_i^{\vee}$, and show that composing $\mathcal N$ with the isomorphism of part (b) gives the pairing of Problem~\ref{s:pairing}. What does $\mathcal N$ do when $\dim(V)=1$, $\operatorname{char}(\KK)=p>0$, and $d=p$?
\item Let $\phi:V\to W$ be a linear map between finite-dimensional spaces. Show that $(\phi^{\vee})^{\otimes d}$ restricts to a map $\Gamma^d(W^{\vee})\to\Gamma^d(V^{\vee})$, and verify that the isomorphisms of part (b) intertwine this map with $\Sym^d(\phi)^{\vee}$. Explain why no division by $d!$ is needed in this statement.
\end{enumerate}
\end{problems}

\hrulefill

\end{document}
